% Relatório técnico-científico (ABNT NBR 10719) com abnTeX2.
% Compilar a partir de docs/:  latexmk -pdf relatorio.tex
% A declaração de uso de IA (recursos/declaracao_ia.tex e recursos/registro_ia.tex)
% é gerada a partir de registro_ia.toml por:  python3 docs/gerar_declaracao_ia.py

\documentclass[
    12pt,
    oneside,
    a4paper,
    openany,
    chapter=TITLE,
    english,
    brazil
]{abntex2}

% ─── Pacotes ────────────────────────────────────────
\usepackage{cmap}
\usepackage[utf8]{inputenc}
\usepackage[T1]{fontenc}
\usepackage{lmodern}
\usepackage{indentfirst}
\usepackage{xcolor}
\usepackage{graphicx}
% protrusion=false: sem pontuação avançando na margem, para respeitar as margens da NBR 14724 à risca
\usepackage[protrusion=false]{microtype}
\usepackage{booktabs}
\usepackage{siunitx}
\usepackage{amsmath}

% Ambiente "quadro" (informação textual com bordas fechadas, NBR 14724), definido
% pelo \newfloat da classe memoir antes de carregar o float, para herdar o
% estilo de legenda do abnTeX2 (título acima, separado por travessão)
% Os quadros entram na Lista de ilustrações (lof), junto com as figuras
\newfloat[chapter]{quadro}{lof}{Quadro}
\renewcommand{\thequadro}{\arabic{quadro}}
\providecommand{\quadroautorefname}{Quadro}
% O babel brazil usa "seção" minúsculo no \autoref; maiúsculo como Capítulo, Tabela etc.
\addto\captionsbrazil{%
    \renewcommand{\sectionautorefname}{Seção}%
    \renewcommand{\subsectionautorefname}{Seção}%
}
\newlistentry{quadro}{lof}{0}
\renewcommand{\cftquadroname}{Quadro\space}
\renewcommand*{\cftquadroaftersnum}{\hfill\textendash\hfill}
\setlength{\cftquadroindent}{\cftfigureindent}
\setlength{\cftquadronumwidth}{\cftfigurenumwidth}

\usepackage{float}
\usepackage{listings}
% Citações conforme a NBR 10520:2023: sobrenome só com inicial maiúscula também
% entre parênteses, ex.: (Cormen et al., 2012), e "et al." em itálico
% (a opção nbr10520-2023 está declarada no início do referencias.bib)
\usepackage[alf, nbr10520-2023, abnt-etal-text=it]{abntex2cite}
% NBR 6023:2018: URL sem os sinais < >
\def\UrlLeft{}
\def\UrlRight{}
% Autor-entidade: maiúsculas e minúsculas na citação, maiúsculas nas referências
% (o \emreferenciastrue fica logo antes do \bibliography; o comando é robusto
% para chegar sem expandir ao .aux e só decidir a caixa na hora de imprimir)
\newif\ifemreferencias
\DeclareRobustCommand{\entidade}[1]{\ifemreferencias\MakeUppercase{#1}\else#1\fi}

\sisetup{
    output-decimal-marker = {,},
    group-digits          = integer,
    group-separator       = {.},
    group-minimum-digits  = 4,
    exponent-product      = \times,
    table-number-alignment = center,
}

\lstset{
    basicstyle   = \ttfamily\footnotesize,
    frame        = single,
    breaklines   = true,
    columns      = fullflexible,
    keepspaces   = true,
    showstringspaces = false,
    % a moldura é desenhada fora da caixa do código; recua o código para ela caber na margem
    xleftmargin  = \dimexpr\fboxsep+\fboxrule\relax,
    xrightmargin = \dimexpr\fboxsep+\fboxrule\relax,
    inputencoding = utf8,
    extendedchars = true,
    literate = {á}{{\'a}}1 {ã}{{\~a}}1 {ç}{{\c c}}1 {é}{{\'e}}1 {ê}{{\^e}}1
               {í}{{\'i}}1 {ó}{{\'o}}1 {õ}{{\~o}}1 {ú}{{\'u}}1 {×}{{$\times$}}1,
}

% Cabeçalho de tabela em mais de uma linha, alinhado pela base: \cabecalho{Desvio\\padrão}
\newcommand{\cabecalho}[1]{\begin{tabular}[b]{@{}c@{}}#1\end{tabular}}

% Marca itens que ainda precisam ser completados antes da entrega
\newcommand{\pendente}[1]{\textcolor{red}{\textbf{[PENDENTE: #1]}}}

% Texto literal de prompts e saídas no apêndice de uso de IA (gerado por
% gerar_declaracao_ia.py, já escapado): recuado, em fonte menor e sem recuo de parágrafo
\newenvironment{registroia}
    {\begin{quote}\small\setlength{\parindent}{0pt}\raggedright}
    {\end{quote}}

% ─── Dados do trabalho ──────────────────────────────
\titulo{Análise empírica das árvores AVL e Rubro-Negra: tempos de construção e busca, altura, rotações e regressão}
\autor{Sérgio de Paiva Ximenes Filho}
\local{Fortaleza}
\data{2026}
% PENDENTE: confirmar orientador, disciplina e atividade
\orientador{Prof. Bruno Lopes}
\instituicao{%
    Universidade de Fortaleza -- UNIFOR
    \par
    Centro de Ciências Tecnológicas
    \par
    Bacharelado em Ciência da Computação}
\tipotrabalho{Relatório técnico}
\preambulo{Relatório da atividade prática de implementação apresentado à disciplina T198 -- Construção e Análise de Algoritmos, do curso de Bacharelado em Ciência da Computação da Universidade de Fortaleza.}

\usepackage{hyperref}
\hypersetup{
    pdftitle    = {Análise empírica das árvores AVL e Rubro-Negra},
    pdfsubject  = {Relatório técnico -- T198 Construção e Análise de Algoritmos},
    colorlinks  = true,
    linkcolor   = black,
    citecolor   = black,
    urlcolor    = black,
    bookmarksdepth = 4,
}

% NBR 14724: margens de 3 cm (sup.) e 2 cm (inf.). O \checkandfixthelayout padrão
% arredonda a altura do texto para um número inteiro de linhas, para cima, e a
% margem inferior caía para 1,82 cm; [fixed] mantém a altura exata
\setulmarginsandblock{3cm}{2cm}{*}
\checkandfixthelayout[fixed]

% NBR 14724: só o número da página, no canto superior direito
\makeevenhead{abntheadings}{}{}{\ABNTEXfontereduzida\thepage}
\makeoddhead{abntheadings}{}{}{\ABNTEXfontereduzida\thepage}
\makeheadrule{abntheadings}{\textwidth}{0pt}

\setlength{\parindent}{1.25cm}
\setlength{\parskip}{0.2cm}
% Folga extra de espaçamento, usada só quando uma linha não cabe na margem
% (ex.: caminhos em \texttt, que não esticam)
\setlength{\emergencystretch}{2em}

% ═══════════════════════════════════════════════════
\begin{document}
\selectlanguage{brazil}
\frenchspacing

% ─── PRÉ-TEXTUAIS ──────────────────────────────────
\pretextual

\imprimircapa
\imprimirfolhaderosto

\setlength{\absparsep}{18pt}
\begin{resumo}
\pendente{resumo: objetivo, método (árvores, cenários, tamanhos, repetições, medições), principais resultados numéricos e conclusão, em um único parágrafo}

\textbf{Palavras-chave}: árvores binárias de busca balanceadas; árvore AVL; árvore rubro-negra; análise empírica; regressão.
\end{resumo}

\pdfbookmark[0]{\listfigurename}{lof}
\listoffigures*
\cleardoublepage

\pdfbookmark[0]{\listtablename}{lot}
\listoftables*
\cleardoublepage

\pdfbookmark[0]{\contentsname}{toc}
\tableofcontents*
\cleardoublepage

% ─── TEXTUAIS ──────────────────────────────────────
\textual

% ═══════════════════════════════════════════════════
\chapter{Introdução}
\label{cap:introducao}

\pendente{contexto: árvores binárias de busca e o problema da degeneração; balanceamento como garantia de altura $O(\log n)$; AVL \cite{adelson1962} e Rubro-Negra \cite{bayer1972,guibas1978,cormen2012}}

\pendente{objetivo e hipóteses esperadas pela teoria (altura, rotações, tempos de construção e busca nos três cenários)}

O \autoref{cap:fundamentacao} resume a teoria das duas árvores. O \autoref{cap:metodologia} descreve o ambiente e o procedimento experimental. O \autoref{cap:resultados} apresenta os dados coletados, sem interpretação. O \autoref{cap:analise} interpreta os resultados à luz da teoria, e o \autoref{cap:conclusao} sintetiza as conclusões.

% ═══════════════════════════════════════════════════
\chapter{Fundamentação teórica}
\label{cap:fundamentacao}

\section{Árvore AVL}
\label{sec:avl}

\pendente{definição, fator de balanceamento, rotações simples e duplas, limite de altura $h < 1{,}44 \log_2(n + 2)$}

\section{Árvore Rubro-Negra}
\label{sec:rb}

\pendente{as cinco propriedades, altura negra, limite $h \le 2 \log_2(n + 1)$, número de rotações por inserção e remoção}

% ═══════════════════════════════════════════════════
\chapter{Metodologia}
\label{cap:metodologia}

Este capítulo descreve o ambiente de execução, a implementação, a geração das entradas, a coleta dos dados e os métodos estatísticos empregados, com o detalhamento necessário para que o experimento possa ser reproduzido. O \autoref{apendice:reproducao} lista os comandos utilizados.

\section{Ambiente de execução}

\pendente{confirmar se o ambiente é o mesmo do Sorting-benchmark e atualizar o quadro}

\begin{quadro}[htbp]
    \caption{Especificações do ambiente de execução}
    \label{quadro:ambiente}
    \centering
    \small
    \begin{tabular}{|l|p{9.5cm}|}
        \hline
        \textbf{Item} & \textbf{Especificação} \\ \hline
        Processador & Intel Core i5-12400 (12.ª geração), frequência base de 2,5 GHz, 6 núcleos e 12 \textit{threads}, cache L2 de 7,5 MiB e L3 de 18 MiB \\ \hline
        Memória RAM & 32 GB instalados (31,8 GB utilizáveis pelo Windows), dos quais 15,5 GiB disponíveis para a máquina virtual do WSL 2 \\ \hline
        Sistema hospedeiro & Microsoft Windows 11 (versão 10.0.26200) \\ \hline
        Sistema de execução & Ubuntu 26.04.1 LTS no Windows Subsystem for Linux 2 (WSL 2), núcleo Linux 6.18.33.2-microsoft-standard-WSL2 \\ \hline
        Compilador & GCC 15.2.0 \\ \hline
        Flags de compilação & \texttt{-Wall -Wextra -O2} \\ \hline
        Análise de dados & Python 3.14.4 com NumPy 2.5.3, pandas 3.0.6, SciPy 1.18.1 e Matplotlib 3.11.2 \\ \hline
    \end{tabular}
    \fonte{Elaborado pelo autor (2026).}
\end{quadro}

O programa foi compilado pelo \texttt{Makefile} do projeto com o comando \texttt{gcc -Wall -Wextra -O2 main.c avl.c rb.c -o main}. \pendente{justificativa do \texttt{-O2} \cite{gcc2025} e condições de uso da máquina durante as execuções}

\section{Implementação das árvores}

\pendente{\texttt{avl.c} e \texttt{rb.c}: estrutura dos nós, inserção, busca, remoção, cálculo da altura (em arestas), contagem de rotações (dupla conta como duas) e sentinela \texttt{nil} da Rubro-Negra \cite{cormen2012}}

\section{Validação}
\label{sec:validacao}

\pendente{o que \texttt{./main validar} verifica em cada árvore e o teste negativo}

\section{Medição}

\pendente{medição em lote com \texttt{clock\_gettime(CLOCK\_MONOTONIC)} \cite{manclock}; o que fica dentro e fora do intervalo medido; construção, $n$ buscas e remoção}

\section{Geração das entradas}
\label{sec:entradas}

\pendente{cenários crescente, decrescente e aleatório (Fisher-Yates); \texttt{srand(repeticao)}; mesmas chaves, buscas e ordem de remoção para as duas árvores}

\section{Coleta e organização dos dados}

\pendente{tamanhos de \num{1000} a \num{1000000}, 50 repetições, ordem do laço, datas da coleta, formato de \texttt{resultados/<arvore>/<cenario>/<n>.csv}}

\section{Estatística descritiva e regressão}
\label{sec:estatistica}

\pendente{média, desvio padrão, IC de 95\,\% com $t$ de Student \cite{montgomery2016}, teste de Welch \cite{welch1947}; regressões do tempo ($n \log_2 n$, $n$, $n^2$) e da altura ($\log_2 n$) com pandas, NumPy, SciPy e Matplotlib \cite{mckinney2010,harris2020,virtanen2020,hunter2007}}

\section{Uso de inteligência artificial}
\label{sec:uso-ia}

\IfFileExists{recursos/declaracao_ia.tex}
    {% Gerado por docs/gerar_declaracao_ia.py a partir de docs/registro_ia.toml. Não edite à mão.
Neste trabalho foi utilizada a ferramenta de inteligência artificial generativa Claude Code (Claude Opus 5.5, claude-opus-5-5), da Anthropic, em 2 sessões e 2 interações, nas datas 23/11/2026. O \autoref{quadro:uso-ia} resume a finalidade de cada interação e a verificação feita pelo autor, e o \autoref{apendice:registro-ia} reproduz a prompt literal e a saída de cada uma. Todo o conteúdo produzido com auxílio da ferramenta foi revisado pelo autor, que se responsabiliza integralmente pelo trabalho.

\textbf{Sessão S1} -- interface: \pendente{informar a interface usada (ex.: claude.ai/code na web, Claude Code no terminal ou no VS Code) e a versão}. Configurações: \pendente{busca na web, memória/histórico de chats, modo de permissões, plugins}. Instruções de sistema ou de projeto: \pendente{instruções de projeto ou personalizadas que estavam ativas, se houver (o prompt de sistema padrão da ferramenta não é reproduzido)}.

\textbf{Sessão S2} -- interface: Claude Code 2.1.330, extensão do Visual Studio Code, executando no Ubuntu 26.04.1 LTS (WSL 2). Configurações: modo de permissões automático; ferramenta de busca na web disponível, mas não usada; memória persistente do Claude Code ativada (vazia no início da sessão); plugin academico-br habilitado; temperatura e demais parâmetros de amostragem no padrão da ferramenta (não configuráveis pelo usuário). Instruções de sistema ou de projeto: prompt de sistema padrão do Claude Code (não reproduzido). Nenhum CLAUDE.md no projeto no início da sessão; o CLAUDE.md criado nesta sessão passa a valer para as seguintes.

\begin{quadro}[htbp]
    \caption{Interações com a ferramenta de inteligência artificial}
    \label{quadro:uso-ia}
    \centering
    \footnotesize
    \begin{tabular}{|c|c|c|p{4.3cm}|p{5.8cm}|}
        \hline
        \textbf{Nº} & \textbf{Sessão} & \textbf{Data} & \textbf{Finalidade} & \textbf{Verificação humana} \\ \hline
        I1 & S1 & 23/11/2026 & Geração do código: estrutura inicial do repositório (árvores AVL e Rubro-Negra em C, benchmark, validação e análise em Python) & \pendente{descrever o que foi conferido (ex.: leitura do código, execução de ./main validar, comparação com o CLRS), corrigido ou descartado.} \\ \hline
        I2 & S2 & 23/11/2026 & Estrutura do relatório em LaTeX (abnTeX2) e criação deste registro de uso de IA & \pendente{descrever o que foi conferido, corrigido ou descartado.} \\ \hline
    \end{tabular}
    \fonte{Elaborado pelo autor (2026).}
\end{quadro}
}
    {\pendente{rodar \texttt{python3 docs/gerar\_declaracao\_ia.py}}}

% ═══════════════════════════════════════════════════
\chapter{Resultados}
\label{cap:resultados}

\pendente{tabelas de estatísticas e gráficos comparativos de tempo de construção, tempo de busca, altura e rotações (e remoção, bônus) por cenário}

% ═══════════════════════════════════════════════════
\chapter{Análise}
\label{cap:analise}

\section{Altura frente aos limites teóricos}
\pendente{}

\section{Rotações}
\pendente{}

\section{Tempos de construção e busca}
\pendente{}

\section{Diferenças entre os cenários}
\pendente{}

\section{Variabilidade e intervalos de confiança}
\pendente{}

\section{Regressão}
\pendente{}

% ═══════════════════════════════════════════════════
\chapter{Conclusão}
\label{cap:conclusao}

\pendente{síntese dos resultados, limitações e trabalhos futuros}

% ─── PÓS-TEXTUAIS ──────────────────────────────────
\postextual

% Entradas citadas só nos apêndices gerados ou ainda não citadas no texto
\nocite{knuth1998,johnson2002,mcgeoch2012,anthropic2026}

\emreferenciastrue
\bibliography{referencias}
\emreferenciasfalse

\begin{apendicesenv}
\partapendices

\chapter{Reprodução do experimento}
\label{apendice:reproducao}

Os comandos a seguir, executados na raiz do projeto, reproduzem o ciclo completo de compilação, validação, execução, análise e geração deste relatório.

\begin{lstlisting}[language=bash]
# 1. Compilar, validar e executar o benchmark (grava os CSVs em resultados/)
make
./main validar
./main

# 2. Preparar o ambiente Python e gerar tabelas e figuras
python3 -m venv venv
venv/bin/pip install -r requirements.txt
venv/bin/python analise.py                # saídas em saida/

# 3. Gerar a declaração de uso de IA e compilar o relatório
python3 docs/gerar_declaracao_ia.py       # saídas em docs/recursos/
cd docs && latexmk -pdf relatorio.tex
\end{lstlisting}

\chapter{Registro do uso de inteligência artificial}
\label{apendice:registro-ia}

\IfFileExists{recursos/registro_ia.tex}
    {% Gerado por docs/gerar_declaracao_ia.py a partir de docs/registro_ia.toml. Não edite à mão.
Este apêndice reproduz, na ordem em que foram enviadas, todas as prompts direcionadas ao desenvolvimento do projeto, com a saída da ferramenta (completa ou a parte usada), a finalidade e a verificação feita pelo autor. O registro original, em formato legível por máquina, está no arquivo \texttt{docs/registro\_ia.toml} do repositório.

\section*{Interação I1 (sessão S1, 23/11/2026)}
\addcontentsline{toc}{section}{Interação I1}
{\raggedright
\begin{description}
    \item[Ferramenta:] Claude Code, Claude Opus 5.5 (\texttt{claude-opus-5-5}), Anthropic
    \item[Interface:] \pendente{informar a interface usada (ex.: claude.ai/code na web, Claude Code no terminal ou no VS Code) e a versão}
    \item[Finalidade:] Geração do código: estrutura inicial do repositório (árvores AVL e Rubro-Negra em C, benchmark, validação e análise em Python)
    \item[Arquivos:] \texttt{main.c}, \texttt{avl.c}, \texttt{avl.h}, \texttt{rb.c}, \texttt{rb.h}, \texttt{analise.py}, \texttt{Makefile}, \texttt{requirements.txt}, \texttt{README.md}, \texttt{.gitignore}
    \item[Commits:] \texttt{fcec72d}
\end{description}
}
\textbf{Prompt literal:}
\begin{registroia}
\pendente{colar aqui a prompt literal que gerou o repositório (a transcrição dessa sessão não está nesta máquina).}
\end{registroia}
\textbf{Saída:}
\begin{registroia}
Commit fcec72d ("Estrutura inicial: AVL e Rubro-Negra com benchmark e análise"), com os arquivos main.c, avl.c, avl.h, rb.c, rb.h, analise.py, Makefile, requirements.txt, README.md e .gitignore. O código completo gerado está no histórico do repositório.
\end{registroia}
\textbf{Verificação humana:} \pendente{descrever o que foi conferido (ex.: leitura do código, execução de ./main validar, comparação com o CLRS), corrigido ou descartado.}

\section*{Interação I2 (sessão S2, 23/11/2026)}
\addcontentsline{toc}{section}{Interação I2}
{\raggedright
\begin{description}
    \item[Ferramenta:] Claude Code, Claude Opus 5.5 (\texttt{claude-opus-5-5}), Anthropic
    \item[Interface:] Claude Code 2.1.330, extensão do Visual Studio Code, executando no Ubuntu 26.04.1 LTS (WSL 2)
    \item[Finalidade:] Estrutura do relatório em LaTeX (abnTeX2) e criação deste registro de uso de IA
    \item[Arquivos:] \texttt{docs/relatorio.tex}, \texttt{docs/referencias.bib}, \texttt{docs/abntex2-alf.bst}, \texttt{docs/registro\_ia.toml}, \texttt{docs/gerar\_declaracao\_ia.py}, \texttt{CLAUDE.md}
\end{description}
}
\textbf{Prompt literal:}
\begin{registroia}
Na Tree-benchmark, crie a pasta docs, vamos usar o latex da mesma forma que foi usado no Sorting Benchmark, nisso, adicione tambem um arquivo para guardar todas as prompts e metados afim de gerar uma declaração de uso de IA no relatório usando esse modelo aqui.

1. O que registrar a cada uso \newline
Item\quad{}Exemplo \newline
Ferramenta e versão exata\quad{}Claude Sonnet 5.5 (claude-sonnet-5-5) \newline
Empresa e interface\quad{}Anthropic, claude.ai (web) ou API \newline
Data(s) de uso\quad{}23/11/2026 \newline
Configurações relevantes\quad{}Busca na web ligada ou desligada, memória ou histórico de chats ligado, parâmetros (temperatura) se usou a API \newline
Prompt literal\quad{}Texto completo, incluindo instruções de sistema ou de projeto, se houver \newline
Saída\quad{}Resposta completa, ou a parte usada \newline
Finalidade\quad{}Revisão de linguagem, síntese de literatura, código, análise de dados etc. \newline
Verificação humana\quad{}O que você checou, corrigiu ou descartou

Toda e qualquer prompt direcionada para o desenvolvimento do projeto precisa ser registrada a partir de agora, isso inclui a prompt que gerou o repositório.
\end{registroia}
\textbf{Saída:}
\begin{registroia}
Criou a pasta docs/ com: relatorio.tex (esqueleto do relatório técnico em abnTeX2, com o mesmo preâmbulo, formatação ABNT e organização em capítulos do relatório do Sorting-benchmark, conteúdo marcado como pendente); referencias.bib e abntex2-alf.bst (copiados e adaptados do Sorting-benchmark, com referências sobre AVL e Rubro-Negra); registro\_ia.toml (este registro); gerar\_declaracao\_ia.py (gera, a partir do registro, a seção de declaração de uso de IA e o apêndice com as prompts literais). Criou também o CLAUDE.md na raiz, instruindo as próximas sessões a registrarem cada prompt aqui.
\end{registroia}
\textbf{Verificação humana:} \pendente{descrever o que foi conferido, corrigido ou descartado.}
}
    {\pendente{rodar \texttt{python3 docs/gerar\_declaracao\_ia.py}}}

\end{apendicesenv}

\end{document}
